
\begin{noteblock}
本目录收录书中正文与附录中提及的外部本体资源线索，供读者查阅与评估；它不是
Semantica package registry，也不是“链接存在即可导入”的自动信任清单。
\end{noteblock}

\section{与 Semantica 的边界}

本书相关的本体、CQ、形状、查询、案例和规则，其唯一可执行正本位于 Semantica。
这里列出的 DOLCE、BFO、SUMO、IOF 等属于候选复用来源。将任一外部资源用于工程包前，
必须核对版本、许可、维护状态、命名空间、语义 profile、来源哈希、与现有概念的冲突，
再通过受治理的 package delta、回归和 release verdict 纳入；不得把下载到书目录的文件
当成平行运行时，也不得因“可被 Protégé 打开”就声称通过 Semantica 验证。


\section{一、通用顶层本体}

\subsection{7. DOLCE（语言学与认知工程的描述性本体）}

\begin{center}
\begin{tabularx}{\linewidth}{@{}lX@{}}
\toprule
\textbf{项目} & \textbf{内容} \\
\midrule
全称 & Descriptive Ontology for Linguistic and Cognitive Engineering \\
开发机构 & 意大利国家研究委员会 ISTC-CNR 应用本体实验室 \\
理论基础 & 受认知科学与语言学启发 \\
核心概念 & 对象、事件、属性、时间、空间等通用基础概念 \\
资源链接 & \url{https://www.loa.istc.cnr.it/index.php/dolce/} \\
\bottomrule
\end{tabularx}
\end{center}
\color{black}

\textbf{适用场景：} 适配所有工程领域本体的顶层概念复用。提供严格的逻辑定义，可有效保障本体的一致性与严谨性，可作为各类工程本体的顶层框架。


\subsection{8. BFO（基础形式本体）}

\begin{center}
\begin{tabular}{@{}ll@{}}
\toprule
\textbf{项目} & \textbf{内容} \\
\midrule
全称 & Basic Formal Ontology \\
标准化 & ISO/IEC 21838-2（2021 年） \\
核心概念 & 实体、过程、属性的基础分类 \\
设计理念 & 结构简洁、逻辑严密 \\
资源链接 & \url{https://basic-formal-ontology.org/} \\
\bottomrule
\end{tabular}
\end{center}
\color{black}

\textbf{适用场景：} 适配机械工程、电力工程、土木工程等各类工程领域。与 OWL 语言完美兼容，可直接作为工程本体的顶层概念架构，减少顶层设计工作量。被多个领域本体（包括 IOF、OBO Foundry 成员本体）作为顶层框架复用。


\subsection{9. SUMO（建议上层合并本体）}

\begin{center}
\begin{tabular}{@{}ll@{}}
\toprule
\textbf{项目} & \textbf{内容} \\
\midrule
全称 & Suggested Upper Merged Ontology \\
规模 & 超过 1 000 个类、2 000 个属性 \\
覆盖领域 & 哲学、数学、计算机、工程等多个领域通用概念 \\
资源链接 & \url{http://www.ontologyportal.org/} \\
\bottomrule
\end{tabular}
\end{center}
\color{black}

\textbf{适用场景：} 适配大型工程本体项目的顶层概念复用，支持复杂逻辑推理，可与多个领域本体集成，适合需要跨领域知识融合的工程场景。


\section{二、领域本体资源库}

\subsection{1. BioPortal（生物医学领域）}

\begin{center}
\begin{tabular}{@{}ll@{}}
\toprule
\textbf{项目} & \textbf{内容} \\
\midrule
定位 & 全球最大生物医学本体库 \\
收录规模 & 超过 1 500 个本体资源（截至 2025 年） \\
覆盖领域 & 临床医学、生命科学、公共卫生等 \\
核心功能 & 本体检索、下载、在线浏览与可视化 \\
资源链接 & \url{https://bioportal.bioontology.org/} \\
\bottomrule
\end{tabular}
\end{center}
\color{black}

\textbf{适用场景：} 适配工程领域中涉及生物医学交叉的场景（如医疗设备工程、生物制药工程）。所有本体均提供详细文档与使用说明，支持商用复用（需遵循对应许可协议）。


\subsection{4. OBO Foundry（生物医学开放本体）}

\begin{center}
\begin{tabular}{@{}ll@{}}
\toprule
\textbf{项目} & \textbf{内容} \\
\midrule
定位 & 生物医学领域开放本体联盟 \\
质量标准 & 本体均遵循严格质量标准 \\
覆盖领域 & 解剖学、细胞学、病理学、药物学等 \\
授权方式 & 全部开源免费，支持二次开发与商用复用 \\
资源链接 & \url{https://obofoundry.org/} \\
\bottomrule
\end{tabular}
\end{center}
\color{black}

\textbf{适用场景：} 适配医疗工程、生物工程等交叉领域的高精度本体需求，提供详细的本体文档与维护更新记录。


\subsection{10. Industrial Ontologies Foundry（IOF，工业本体基础）}

\begin{center}
\begin{tabular}{@{}ll@{}}
\toprule
\textbf{项目} & \textbf{内容} \\
\midrule
定位 & 面向制造业、供应链、工业工程的开放本体联盟 \\
建立模式 & 参照 OBO Foundry 模式，核心本体基于 BFO 顶层框架 \\
核心概念 & 工业过程、设备、资源、能力等核心工程概念 \\
维护方式 & 工业界与学术界联合维护 \\
资源链接 & \url{https://www.industrialontologies.org/} \\
\bottomrule
\end{tabular}
\end{center}
\color{black}

\textbf{适用场景：} 适配智能制造、数字孪生、工业互联网等场景。是目前工程本体领域最具行业影响力的开放标准资源，可直接复用或作为自建工程本体的参考框架。


\subsection{11. IEEE 1872 机器人与自动化标准本体}

\begin{center}
\begin{tabularx}{\linewidth}{@{}lX@{}}
\toprule
\textbf{项目} & \textbf{内容} \\
\midrule
发布机构 & IEEE 标准协会 \\
标准版本 & 1872-2015（首版）、1872.1-2024（任务知识表示）、1872.2-2021（自主机器人本体） \\
覆盖领域 & 机器人学与自动化领域核心概念、关系与公理体系 \\
特点 & 具有强行业权威性，便于与行业系统互操作 \\
\bottomrule
\end{tabularx}
\end{center}
\color{black}

\textbf{资源链接：}

\begin{center}
\begin{tabularx}{\linewidth}{@{}XX@{}}
\toprule
\textbf{资源} & \textbf{链接} \\
\midrule
IEEE 1872-2015 官方标准 & \url{https://standards.ieee.org/standard/1872-2015.html} \\
IEEE 1872.1-2024 官方标准 & \url{https://standards.ieee.org/ieee/1872.1/6993/} \\
开源 OWL 文件（可直接导入 Protégé） & \url{https://opensource.ieee.org/aur/owl} \\
\bottomrule
\end{tabularx}
\end{center}
\color{black}

\textbf{适用场景：} 适配工业机器人、智能制造、自动化系统等工程场景，适合需要与自动化标准对接的工程本体项目。开源 OWL 文件可直接下载后导入 Protégé 等本体工具使用。


\section{三、通用词汇表与 Web 本体}

\subsection{2. Linked Open Vocabularies（LOV，语义网通用词汇表）}

\begin{center}
\begin{tabular}{@{}ll@{}}
\toprule
\textbf{项目} & \textbf{内容} \\
\midrule
定位 & 语义网领域通用 RDF 词汇表目录 \\
收录规模 & 约 600 余个词汇表 \\
覆盖范围 & 通用概念、领域专用词汇 \\
标准合规 & 符合 W3C 标准 \\
资源链接 & \url{https://lov.linkeddata.es/} \\
\bottomrule
\end{tabular}
\end{center}
\color{black}

\textbf{适用场景：} 适配各类工程本体的通用概念复用（如"时间""空间""实体"等基础概念），兼容性强，可直接导入 Protégé 等工具。

\begin{noteblock}
\textbf{注意：} LOV 收录的是轻量级 RDF 词汇表，部分条目并非完整的工程级本体，选用时应结合领域需求评估适配性。
\end{noteblock}


\subsection{3. Schema.org（Web 内容标记本体）}

\begin{center}
\begin{tabular}{@{}ll@{}}
\toprule
\textbf{项目} & \textbf{内容} \\
\midrule
发起方 & 谷歌、微软、雅虎等企业联合推出 \\
定位 & 聚焦 Web 内容标记与数据互通 \\
核心概念 & 产品、服务、组织、事件等通用概念 \\
设计理念 & 结构简洁、易用性强 \\
资源链接 & \url{https://schema.org/} \\
\bottomrule
\end{tabular}
\end{center}
\color{black}

\textbf{适用场景：} 适配工程领域中涉及 Web 展示、数据共享的本体项目（如工程设备展示、项目信息发布），支持与各类 Web 系统集成，无需复杂适配。


\section{四、多领域本体仓库与工具}

\subsection{5. OntoHub（多领域异构本体仓库）}

\begin{center}
\begin{tabular}{@{}ll@{}}
\toprule
\textbf{项目} & \textbf{内容} \\
\midrule
技术特点 & 支持 DOL（分布式本体语言） \\
收录规模 & 超过 2 万个本体 \\
覆盖领域 & 工程、计算机、商业、教育等多个领域 \\
核心功能 & 版本管理、在线编辑、协作开发、检索下载 \\
资源链接 & \url{https://www.ontohub.org/} \\
\bottomrule
\end{tabular}
\end{center}
\color{black}

\textbf{适用场景：} 定位偏向学术研究场景，适合需要异构本体整合、跨语言本体互操作的工程本体项目。支持按领域筛选本体，提供本体质量评估报告，方便快速判断适配性。


\subsection{6. Protégé 本体库（教学与示例本体，已停止维护）}

\begin{center}
\begin{tabular}{@{}ll@{}}
\toprule
\textbf{项目} & \textbf{内容} \\
\midrule
定位 & Protégé 工具官方配套本体库 \\
内容 & 大量教学用示例本体，涵盖基础概念与简单领域本体 \\
资源链接 & \url{https://protege.stanford.edu/ontologies/} \\
维护状态 &  \textbf{已停止维护更新，不建议用于生产项目} \\
\bottomrule
\end{tabular}
\end{center}
\color{black}

\textbf{适用场景：} 适合初学者入门练习。

\begin{noteblock}
\textbf{替代建议：}
\begin{itemize}
\item 查找生物医学相关示例本体 \(\rightarrow\) 转至 \textbf{\href{https://bioportal.bioontology.org/}{BioPortal}}
\item 查找通用工程示例本体 \(\rightarrow\) 参考 \textbf{\href{https://obofoundry.org/}{OBO Foundry}} 或 \textbf{\href{https://www.industrialontologies.org/}{IOF}}
\end{itemize}

Protégé 工具本身持续维护，可从 \url{https://protege.stanford.edu/} 下载使用。
\end{noteblock}


\section{资源速查表}

\begin{center}
\begin{tabularx}{\linewidth}{@{}lllXX@{}}
\toprule
\textbf{\#} & \textbf{资源名称} & \textbf{类型} & \textbf{适用场景关键词} & \textbf{链接} \\
\midrule
1 & BioPortal & 本体库 & 生物医学、医疗设备、生物制药 & \href{https://bioportal.bioontology.org/}{链接} \\
2 & LOV & 词汇表目录 & 语义网、通用基础概念 & \href{https://lov.linkeddata.es/}{链接} \\
3 & Schema.org & Web 本体 & Web 展示、数据共享 & \href{https://schema.org/}{链接} \\
4 & OBO Foundry & 本体联盟 & 医疗工程、生物工程 & \href{https://obofoundry.org/}{链接} \\
5 & OntoHub & 本体仓库 & 异构本体整合、学术研究 & \href{https://www.ontohub.org/}{链接} \\
6 & Protégé 本体库 & 示例本体 & 入门练习（ 已停维护） & \href{https://protege.stanford.edu/ontologies/}{链接} \\
7 & DOLCE & 顶层本体 & 全领域顶层框架 & \href{https://www.loa.istc.cnr.it/index.php/dolce/}{链接} \\
8 & BFO & 顶层本体 & 机械 / 电力 / 土木工程（ISO 标准） & \href{https://basic-formal-ontology.org/}{链接} \\
9 & SUMO & 顶层本体 & 大型跨领域工程本体 & \href{http://www.ontologyportal.org/}{链接} \\
10 & IOF & 工业本体联盟 & 智能制造、数字孪生、工业互联网 & \href{https://www.industrialontologies.org/}{链接} \\
11 & IEEE 1872 & 标准本体 & 机器人、自动化、智能制造 & \href{https://opensource.ieee.org/aur/owl}{链接} \\
\bottomrule
\end{tabularx}
\end{center}
\color{black}


\textit{本目录随书持续更新。如发现链接失效或有新资源建议，欢迎提交 Issue 或 Pull Request。}
